% !TeX root = appendices.tex
\documentclass[../main.tex]{subfiles}
\graphicspath{{\subfix{../images/}}}

\begin{document}
    \begin{appendices} 
    \section{Nonlinear Framework}
    \subsection{Material Model Expressed in \emph{\textbf{S}}}
    Given the Cauchy stress tensor \eqref{neo_hookean_definition} and the definition of the Second Piola-Kircchhof stress tensor,~\eqref{eq:19}
     can be written as,
    \begin{equation}
        \begin{split}
        \mathrm{\boldsymbol{S}} &= \mathrm{J \boldsymbol{F^{-1}}} \left[ \frac{\mu}{J} \left(\boldsymbol{b} - \boldsymbol{I}\right) + \frac{\Lambda}{J}\ln{J}\boldsymbol{I}\right] \boldsymbol{F^{-T}}\\
        &= \mathrm{J}  \mathrm{\left[ \frac{\mu}{J} \left(\boldsymbol{F^{-1}b} - \boldsymbol{F^{-1}}\right) + \frac{\Lambda}{J}\ln{J}\boldsymbol{F^{-1}}\right] \boldsymbol{F^{-T}}}\\
        &= \mathrm{J} \mathrm{\left[ \frac{\mu}{J} \left(\boldsymbol{F^{-1}}\left(\boldsymbol{FF^T}\right) - \boldsymbol{F^{-1}}\right) + \frac{\Lambda}{J}\ln{J}\boldsymbol{F^{-1}}\right]} \boldsymbol{F^{-T}}\\ 
        &= \mathrm{J} \mathrm{\left[ \frac{\mu}{J} \left(\boldsymbol{F^T} - \boldsymbol{F^{-1}} \right) \boldsymbol{F^{-T}} + \frac{\Lambda}{J} \ln{J}\boldsymbol{F^{-1}F^{-T}} \right]}\\ 
        &= \mu \left( \mathrm{\boldsymbol{F^TF^{-T}}} - \left( \mathrm{\boldsymbol{F^TF}} \right)^{-1} \right) + \Lambda\mathrm{\ln{J}\left(\boldsymbol{F^TF}\right)^{-1}}
        \end{split}
    \end{equation}
    
    \begin{equation}\label{S_neohookean}
        \therefore \mathrm{\boldsymbol{S} = \mu\left(\boldsymbol{I} - \boldsymbol{C^{-1}}\right) + \Lambda\ln{J}\boldsymbol{C^{-1}}}
    \end{equation}
    \subsection{Derivation of Compliance Matrix in Reference Configuration}
    As the formulation assumes the existence of a cartesian basis, the right Cauchy-Green tensor and its inverse can be written, 
    \begin{equation} \label{eq:def_C_comp}
        \begin{split}
        {\boldsymbol{C}} &= {C^i_{.j}\boldsymbol{e}_i \otimes \boldsymbol{e}^j}\\
        {\boldsymbol{C}}^{-1} &= {\left(C^{-1}\right)^i_{.j} \boldsymbol{e}_i \otimes \boldsymbol{e}^j}.
        \end{split}
    \end{equation}
Writing~\eqref{S_neohookean} in indicial form using \eqref{eq:def_C_comp} and using the definition~\eqref{compliance_definition}
    \begin{equation} \label{eq:ref_comp_1}
        \begin{split}
            {S^i_{.j}} &= \mu {\left[\delta ^i_{.i} - \left(C^{-1}\right)^i_{.j}\right] + \Lambda \ln{J} \left(C^{-1}\right)^i_{.j}}\\
        {\Rightarrow 2\frac{\partial S^i_{.j}}{\partial C^a_{.b}}} &= 2\mu {\left[0 - \frac{\partial \left(C^{-1}\right)^i_{.j}}{\partial C^a_{.b}} \right] + 2\Lambda\frac{\partial \left(\ln{J}\left(C^{-1}\right)^i_{.j}\right)}{\partial C^a_{.b}}}\\
            &= 2\mu {\left[0 - \frac{\partial \left(C^{-1}\right)^i_{.j}}{\partial C^a_{.b}} \right] + 2\Lambda \left(C^{-1}\right)^i_{.j} \boldsymbol{e}_i \otimes \boldsymbol{e}^j \otimes \frac{\partial \left(\ln{J}\right)}{\partial C^a_{.b}} + 2\Lambda\ln{J} \frac{\partial \left(C^{-1}\right)^i_{.j}}{\partial C^a_{.b}}}\\
            &= {2\left(-\mu + \Lambda\ln{J}\right) \frac{\partial \left(C^{-1}\right)^i_{.j}}{\partial C^a_{.b}} + 2\Lambda\left(C^{-1}\right)^i_{.j} \boldsymbol{e}_i \otimes \boldsymbol{e}^j \otimes \left[\frac{1}{J} \frac{\partial J}{\partial C^a_{.b}}\right]}\\
        \implies \mathbb{C}^{i.a.}_{.j.b} &= 2\left(-\mu + \Lambda\ln{{J}}\right)\left({C}^{-1}\right)^{i}_{.a} \left({C}^{-1}\right)^{b}_{.j} + \Lambda {\left(C^{-1}\right)^i_{.j} \left(C^{-1}\right)^a_{.b}}
        \end{split}
    \end{equation}
    \subsection{Derivation of Compliance Matrix in Current Configuration}
Given the expression of the compliance matrix in reference configuration
\begin{equation}\label{eq:22}
    \mathbb{C} = 2\left(-\mu + \Lambda\ln{J}\right)\left(C^{-1}\right)^i_{.a} \left(C^{-1}\right)^b_{.j} \boldsymbol{e}_{i} \otimes \boldsymbol{e}^{a} \otimes \boldsymbol{e}_{b} \otimes \boldsymbol{e}^{j}
    + \Lambda \left(C^{-1}\right)^i_{.j} \left(C^{-1}\right)^a_{.b} \boldsymbol{e}_{i} \otimes \boldsymbol{e}^{j} \otimes \boldsymbol{e}_{a} \otimes \boldsymbol{e}^{b}.
\end{equation}
and using the push-forward transformation to transform \eqref{eq:22} to the current configuration
\begin{equation}
    \begin{split}
    \mathbbmss{c} &= 2\left(-\mu + \Lambda\ln{J}\right)\left(C^{-1}\right)^i_{.a} \left(C^{-1}\right)^b_{.j} 
    \left(\boldsymbol{Fe}_i \otimes \boldsymbol{e}^a \boldsymbol{F^T}\right) \otimes \left(\boldsymbol{Fe}_b \otimes \boldsymbol{e}^j \boldsymbol{F^T}\right)
    + \Lambda \left(C^{-1}\right)^i_{.j} \left(C^{-1}\right)^a_{.b} 
    \left(\boldsymbol{Fe}_i \otimes \boldsymbol{e}^j \boldsymbol{F^T}\right) \otimes \left(\boldsymbol{Fe}_a \otimes \boldsymbol{e}^b \boldsymbol{F^T}\right)\\
    &:= \mathbb{I} + \mathbb{II}
    \end{split}
\end{equation}
Operating on $\mathbb{I}$
\begin{equation} \label{eq:I}
    \begin{split}
        \mathbb{I} &= 2\left(-\mu + \Lambda\ln{J}\right)\left(C^{-1}\right)^i_{.a} \left(C^{-1}\right)^b_{.j} \left[F^p_{.q} \boldsymbol{e}_p (\boldsymbol{e}^q \cdot \boldsymbol{e}_i)\right] \otimes \left[(\boldsymbol{e}^a \cdot \boldsymbol{e}_r) \boldsymbol{e}^s \left(F^T\right)^r_{.s}\right]
         \left[F^t_{.u} \boldsymbol{e}_t (\boldsymbol{e}^u \cdot \boldsymbol{e}_b)\right] \otimes \left[(\boldsymbol{e}^j \cdot \boldsymbol{e}_v) \boldsymbol{e}^w \left(F^T\right)^v_{.w} \right]\\
        &= 2\left(-\mu + \Lambda\ln{J}\right) \left(F^{-1}F^{-T}\right)^i_{.a} \left(F^{-1}F^{-T}\right)^b_{.j} \left[F^p_{.q} \delta^q_{.i} \delta^a_{.r} (F^T)^r_{.s}   F^t_{.u} \delta^u_{.b} \delta^j_{.v} (F^T)^v_{.w} \right] \boldsymbol{e}_p \otimes \boldsymbol{e}^s \otimes \boldsymbol{e}_t \otimes \boldsymbol{e}^w\\
        & = 2\left(-\mu + \Lambda\ln{J}\right) \left(F^{-1}\right)^i_{.c} \left(F^{-T}\right)^c_{.a} \left(F^{-1}\right)^b_{.d} \left(F^{-T}\right)^d_{.j} \left[F^p_{.i} \delta^a_{.r} (F^T)^r_{.s} F^t_{.b} \delta^j_{.v} (F^T)^v_{.w}\right] \boldsymbol{e}_p \otimes \boldsymbol{e}^s \otimes \boldsymbol{e}_t \otimes \boldsymbol{e}^w\\
        & = 2\left(-\mu + \Lambda\ln{J}\right)\left[ F^p_{.i} (F^{-1})^i_{.c} F^t_{.b} (F^{-1})^b_{.d}\right] \left[(F^{-T})^c_{.a} \delta^a_{.r} (F^T)^r_{.s}\right] \left[(F^{-T})^d_{.j} \delta^j_{.v} (F^T)^v_{.w}\right] \boldsymbol{e}_p \otimes \boldsymbol{e}^s \otimes \boldsymbol{e}_t \otimes \boldsymbol{e}^w\\
        & =  2\left(-\mu + \Lambda\ln{J}\right) \delta^p_{.c} \delta^t_{.d} (F^{-T})^c_{.r} (F^T)^r_{.s} (F^{-T})^d_{.v} (F^T)^v_{.w} \boldsymbol{e}_p \otimes \boldsymbol{e}^s \otimes \boldsymbol{e}_t \otimes \boldsymbol{e}^w\\
        &= 2\left(-\mu + \Lambda\ln{J}\right) \delta^p_{.c} \delta^t_{.d} \delta^c_{.s} \delta^d_{.w} \boldsymbol{e}_p \otimes \boldsymbol{e}^s \otimes \boldsymbol{e}_t \otimes \boldsymbol{e}^w\\
    \implies \mathbb{I} &= 2\left(-\mu + \Lambda\ln{J}\right) \delta^p_{.s} \delta^t_{.w} \boldsymbol{e}_p \otimes \boldsymbol{e}^s \otimes \boldsymbol{e}_t \otimes \boldsymbol{e}^w
    \end{split}
\end{equation}
Similarly for $\mathbb{II}$
\begin{equation} \label{eq:II}
    \begin{split}
        \mathbb{II} &= \frac{\Lambda}{2}\left(C^{-1}\right)^i_{.j} \left(C^{-1}\right)^a_{.b} \left[F^p_{.q} \boldsymbol{e}_p (\boldsymbol{e}^q \cdot \boldsymbol{e}_i)\right] \otimes \left[(\boldsymbol{e}^j \cdot \boldsymbol{e}_v) \boldsymbol{e}^w \left(F^T\right)^v_{.w}\right] \otimes
    \left[F^r_{.s} \boldsymbol{e}_r (\boldsymbol{e}^s \cdot \boldsymbol{e}_a)\right] \otimes \left[(\boldsymbol{e}^b \cdot \boldsymbol{e}_t) \boldsymbol{e}^u \left(F^T\right)^t_{.u} \right]\\
    &= \frac{\Lambda}{2} \left(C^{-1}\right)^i_{.j} \left(C^{-1}\right)^a_{.b} \left[F^p_{.q} \delta^q_{.i} \delta^j_{.v} (F^T)^v_{.w} F^r_{.s} \delta^s_{.a} \delta^b_{.t} (F^T)^t_{.u} \right] \boldsymbol{e}_p \otimes \boldsymbol{e}^w \otimes \boldsymbol{e}_r \otimes \boldsymbol{e}^u\\
    &= \frac{\Lambda}{2} \left(F^{-1}F^{-T}\right)^i_{.j} \left(F^{-1}F^{-T}\right)^a_{.b} \left[F^p_{.i} \delta^j_{.v} (F^T)^v_{.w} F^r_{.a} \delta^b_{.t} (F^T)^t_{.u} \right]\boldsymbol{e}_p \otimes \boldsymbol{e}^w \otimes \boldsymbol{e}_r \otimes \boldsymbol{e}^u\\
    &= \frac{\Lambda}{2} \left(F^{-1}\right)^i_{.c} \left(F^{-T}\right)^c_{.j} \left(F^{-1}\right)^a_{.d} \left(F^{-T}\right)^d_{.b} \left[F^p_{.i} \delta^j_{.v} (F^T)^v_{.w} F^r_{.a} \delta^b_{.t} (F^T)^t_{.u} \right]\boldsymbol{e}_p \otimes \boldsymbol{e}^w \otimes \boldsymbol{e}_r \otimes \boldsymbol{e}^u\\
    &= \frac{\Lambda}{2}(F^p_{.i} (F^{-1})^i_{.c})(F^r_{.a}(F^{-1})^a_{.d})\left[\left(F^{-T}\right)^c_{.j} \delta^j_{.v} (F^T)^v_{.w}\right] \left[\left(F^{-T}\right)^d_{.b} \delta^b_{.t} (F^T)^t_{.u}\right]\boldsymbol{e}_p \otimes \boldsymbol{e}^w \otimes \boldsymbol{e}_r \otimes \boldsymbol{e}^u\\
    &= \frac{\Lambda}{2} \delta^p_{.c} \delta^r_{.d} \left[(F^{-T})^c_{.v}(F^T)^v_{.w}\right] \left[(F^{-T})^d_{.t}(F^T)^t_{.u}\right]\boldsymbol{e}_p \otimes \boldsymbol{e}^w \otimes \boldsymbol{e}_r \otimes \boldsymbol{e}^u\\
    &= \frac{\Lambda}{2} \delta^p_{.c} \delta^r_{.d} \delta^c_{.w} \delta^d_{.u} \boldsymbol{e}_p \otimes \boldsymbol{e}^w \otimes \boldsymbol{e}_r \otimes \boldsymbol{e}^u\\
    &= \frac{\Lambda}{2} \underbrace{\delta^p_{.w} \delta^r_{.u} \boldsymbol{e}_p \otimes \boldsymbol{e}^w \otimes \boldsymbol{e}_r \otimes \boldsymbol{e}^u}_{r \rightarrow s, u \rightarrow t}\\
\implies \mathbb{II} &= \frac{\Lambda}{2} \delta^p_{.w} \delta^s_{.t} \boldsymbol{e}_p \otimes \boldsymbol{e}^w \otimes \boldsymbol{e}_s \otimes \boldsymbol{e}^t
\end{split}
\end{equation}
Using \eqref{eq:I} and \eqref{eq:II}, the compliance matrix in the current configuration in indicial form is
\begin{equation}
    \mathbbmss{c}^{p.s.}_{.w.t} = 2\left(-\mu + \Lambda\ln{J}\right) \delta^p_{.s} \delta^t_{.w} + \frac{\Lambda}{2} \delta^p_{.w} \delta^s_{.t}.
\end{equation}



    \section{Shell Formulations}
    \subsection{Curviliniear Coordinates}
    \subsubsection{Proof: \texorpdfstring{$\boldsymbol{a}_{\alpha;\beta} = (\boldsymbol{n} \otimes \boldsymbol{n}) \boldsymbol{a}_{\alpha, \beta}$}{Cov Derivative}}
    Using~\eqref{surface_identity}
    \begin{equation}\label{eq:eq12}
        \left( \boldsymbol{n} \otimes \boldsymbol{n} \right) = \boldsymbol{\tilde{1}} - \boldsymbol{1},
    \end{equation}
    it can be written     
    \begin{equation}\label{appendices:cov_tangent_derivative}
        \begin{split}
            \left( \boldsymbol{n} \otimes \boldsymbol{n} \right) \boldsymbol{a}_{\alpha, \beta} &= \left( \boldsymbol{\tilde{1}} - \boldsymbol{1} \right) \boldsymbol{a}_{\alpha, \beta} \\ &= \boldsymbol{\tilde{1}} \boldsymbol{a}_{\alpha, \beta} - \left( \boldsymbol{a}_{\gamma} \otimes \boldsymbol{a}^{\gamma} \right) \boldsymbol{a}_{\alpha, \beta} \\
            &= \boldsymbol{a}_{\alpha, \beta} - \left( \boldsymbol{a}_{\alpha, \beta} \cdot \boldsymbol{a}^{\gamma}\right) \boldsymbol{a}_{\gamma} \\
            &= \boldsymbol{a}_{\alpha, \beta} - \Gamma_{\alpha \beta}^{\gamma} \boldsymbol{a}_{\gamma}\\
            &= \boldsymbol{n} \cdot \boldsymbol{a}_{\alpha ; \beta}
        \end{split}
    \end{equation}
    \subsubsection{Computation of Components of \texorpdfstring{$b_{\alpha \beta}$}{Curvature tensor}}
    Using the results from the previous section of the Appendix,
    \begin{equation}\label{appendices:curvature_tensor_derivation}
        \begin{split}
            \boldsymbol{n} \cdot \boldsymbol{a}_{\alpha ; \beta} &= \boldsymbol{n} \cdot \left( \boldsymbol{n} \otimes \boldsymbol{n} \right) \boldsymbol{a}_{\alpha, \beta} \\
            &= \left( \boldsymbol{n} \cdot \boldsymbol{n} \right) \left( \boldsymbol{a}_{\alpha, \beta} \cdot \boldsymbol{n} \right) \\
            &= \boldsymbol{1} \left( \boldsymbol{a}_{\alpha, \beta} \cdot \boldsymbol{n} \right) =: b_{\alpha \beta}
        \end{split}
    \end{equation}
        
    \subsection{Weak Form Derivation} \label{appendices:weak_form_derivation}
    \begin{equation}\label{eq:weak_1}
        \int_{\mathcal{S}} \boldsymbol{w} \left[ \boldsymbol{t}_{; \alpha}^{\alpha} + \boldsymbol{f} \right] \, \mathrm{d} a = 0.
    \end{equation}
    Decomposing $\boldsymbol{w}$ into its surface and normal components     
    \begin{equation}\label{eq:weak_2}
        0 = \int_{\mathcal{S}} \left[ w_{\gamma} \boldsymbol{a}^{\gamma} + w \boldsymbol{n} \right] \left[ \sigma_{~~; \alpha}^{\beta \alpha} \boldsymbol{a}_{\beta} + \sigma^{\beta \alpha} b_{\beta \alpha} \boldsymbol{n} + \boldsymbol{f} \right] \, \mathrm{d} a
    \end{equation}
        
        \begin{equation}\label{eq:weak_3}
        0 = \int_{\mathcal{S}} w_{\gamma} \boldsymbol{a}^{\gamma} \left[ \sigma^{\beta \alpha}_{~~; \alpha} \boldsymbol{a}_{\beta} + f^{\alpha} \boldsymbol{a}_{\alpha} \right] \, \mathrm{d} a + \int_{\mathcal{S}} w \boldsymbol{n} \left[ \sigma^{\beta \alpha} b_{\beta \alpha} \boldsymbol{n} + p \boldsymbol{n} \right] \, \mathrm{d} a
        \end{equation}
        
        \begin{equation}\label{eq:weak_4}
        \begin{split}
                0 &= \int_{\mathcal{S}} \left( w_{\gamma} \sigma^{\beta \alpha}_{~~; \alpha} \delta^{\gamma}_{\beta} + f^{\alpha} \delta^{\gamma}_{\alpha} \right) \, \mathrm{d} a + \int_{\mathcal{S}} w \left[ \sigma^{\beta \alpha} b_{\beta \alpha} + p \right] \, \mathrm{d} a \\ &= \int_{\mathcal{S}} \left[ w_{\beta} \sigma^{\beta \alpha}_{~~; \alpha} + f^{\gamma} \right] \, \mathrm{d} a + \int_{\mathcal{S}} \left[ w \sigma^{\beta \alpha} b_{\beta \alpha} +  p \right] \, \mathrm{d} a  \\ &= \int_{\mathcal{S}} \left[ w_{\beta} \sigma^{\beta \alpha}_{~~; \alpha} + f^{\gamma} \right] \, \mathrm{d} a + \int_{\mathcal{S}} \left[ w \sigma^{\beta \alpha} b_{\beta \alpha} + p \right] \, \mathrm{d} a \\ &= \int_{\mathcal{S}} \left[ w_{\gamma} \sigma^{\gamma \alpha}_{~~; \alpha} + w_{\gamma} f^{\gamma} \right] \, \mathrm{d} a + \int_{\mathcal{S}} \left[ w \sigma^{\beta \alpha} b_{\beta \alpha} + p \right] \, \mathrm{d} a
            \end{split}
        \end{equation}
        
        Using divergence theorem
        
        \begin{equation}\label{eq:weak_5}
            \begin{split}
                \int_{\mathcal{S}} w_{\gamma} \sigma^{\gamma \alpha}_{~~; \alpha} \, \mathrm{d} a &= \int_{\mathcal{S}} \left( w_{\gamma} \sigma^{\gamma \alpha} \right)_{; \alpha} \, \mathrm{d} a - \int_{\mathcal{S}} w_{\gamma ; \alpha} \sigma^{\gamma \alpha} \, \mathrm{d} a \\ &= \int_{\partial \mathcal{S}} w_{\gamma} \sigma^{\gamma \alpha} m_{\alpha} \, \mathrm{d} s - \int_{\mathcal{S}} w_{\gamma ; \alpha} \sigma^{\gamma \alpha} \, \mathrm{d} a
            \end{split}
        \end{equation}
        
        using $ \boldsymbol{t} = m_{\alpha} \boldsymbol{t}^{\alpha} = m_{\alpha} \sigma^{\beta \alpha} \boldsymbol{a}_{\beta} $ it is possible to rewrite the surfrace integral of~\eqref{eq:weak_5} as
\begin{equation}\label{eq:weak_6}
    \int_{\partial \mathcal{S}} w_{\gamma} \sigma^{\gamma \alpha} m_{\alpha} \, \mathrm{d} s = \int_{\partial \mathcal{S}} w_{\gamma} \overline{t}^{\gamma} \, \mathrm{d} s
\end{equation}
the weak form can be written as,
    
    \begin{equation}\label{eq:wear_7}
    0 = \int_{\partial \mathcal{S}} w_{\gamma} \overline{t}^{\gamma} \, \mathrm{d} s - \int_{\mathcal{S}} w_{\gamma ; \alpha} \sigma^{\gamma \alpha} \, \mathrm{d} a + \int_{\mathcal{S}} w_{\gamma} f^{\gamma} \, \mathrm{d} a + \int_{\mathcal{S}} w \sigma^{\beta \alpha} b_{\beta \alpha} \, \mathrm{d} a + \int_{\mathcal{S}} w p \, \mathrm{d} s.
    \end{equation}
  %%%---------------------------------------
  % w_{a;b}  
    \subsection{Computation of \texorpdfstring{$w_{\alpha ; \beta}$}{w-a;b}}
    \begin{equation}\label{appendices:w_ab}
        \begin{split}
            w_{\alpha ; \beta} &= \left( \boldsymbol{w} \cdot \boldsymbol{a}_{\alpha} \right)_{; \beta} = \boldsymbol{w}_{; \beta} \cdot \boldsymbol{a}_{\alpha} + \boldsymbol{w} \cdot \boldsymbol{a}_{\alpha ; \beta} \\ &= w_{; \beta} \boldsymbol{a}_{\alpha} + \left( w_{\gamma} \boldsymbol{a}^{\gamma} + w \boldsymbol{n} \right) a_{\alpha ; \beta} \\ &= \boldsymbol{w}_{; \beta} \cdot \boldsymbol{a}_{\alpha} + w_{\gamma} \boldsymbol{a}^{\gamma} \cdot \boldsymbol{a}_{\alpha ; \beta} + w \boldsymbol{n} \cdot \boldsymbol{a}_{\alpha ; \beta} \\ &= \boldsymbol{w}_{; \beta} \cdot \boldsymbol{a}_{\alpha} + w b_{\alpha \beta}
        \end{split}
    \end{equation}
    \subsection{Compliance Matrix} \label{appendices:compliance_matrix_curvy}
With the help of equations~\eqref{eq:NH_strain} and~\eqref{eq:CGtrace}, the Kircchhof stress tensor can be written as
\begin{equation}\label{eq:kirch_NH}
    \begin{split}
        \tau^{\alpha \beta} &= 2 \frac{\partial W}{\partial a_{\alpha \beta}} = 2 \left[ \frac{\Lambda}{4} \left( \frac{\partial J^{2}}{\partial a_{\alpha \beta}} - 2 \frac{\partial \ln{J}}{\partial a_{\alpha \beta}} \right) + \frac{\mu}{2} \left( \frac{\partial I_{1}}{\partial a_{\alpha \beta}} - 2 \frac{\partial \ln{J}}{\partial a_{\alpha \beta}} \right) \right] \\ &= \frac{\Lambda}{2} \left( \frac{\partial J^{2}}{\partial J} \frac{\partial J}{\partial a_{\alpha \beta}} - 2 \frac{\partial \ln{J}}{\partial J} \frac{\partial J}{\partial a_{\alpha \beta}} \right) + \mu \left( \frac{\partial I_{1}}{\partial a_{\alpha \beta}} - 2 \frac{\partial \ln{J}}{\partial J} \frac{\partial J}{\partial a_{\alpha \beta}} \right) 
    \end{split}
\end{equation}
Computing the derivative of J as defined in \eqref{area_something_man} with respect to the metric tensor on the surface,
\begin{equation}\label{eq:dJdaab}
    \frac{\partial J}{\partial a_{\alpha \beta}} = \frac{1}{J_{A}} \frac{\partial J_{a}}{\partial a_{\alpha \beta}} = \frac{1}{J_{A}} \frac{\partial \sqrt{\det{a_{\alpha \beta}}}}{\partial \det{a_{\alpha \beta}}} \frac{\partial \det{a_{\alpha \beta}}}{\partial a_{\alpha \beta}} = \frac{1}{J_{A}} \frac{1}{2} \det{a_{\alpha \beta}}^{-\frac{1}{2}} \det{a_{\alpha \beta}} \left( a_{\alpha \beta}^{-1} \right)^{T} = \frac{1}{2 J_{A}} J_{a} a^{\alpha \beta} = \frac{J}{2} a^{\alpha \beta}
\end{equation}

and using~\eqref{eq:CGtrace} and~\eqref{eq:dJdaab}, equation~\eqref{eq:kirch_NH} can be written as

\begin{equation}\label{eq:kirch_NH_2}
    \begin{split}
        \tau^{\alpha \beta} &= \frac{\Lambda}{2} \left( 2 J \frac{J}{2} a^{\alpha \beta} - \frac{1}{J} \frac{J}{2} a^{\alpha \beta} \right) + \mu \left( A^{\alpha \beta} - 2 \frac{1}{J} \frac{J}{2} a^{\alpha \beta} \right) \\ &= \frac{\Lambda}{2} \left( J^{2} - 1 \right) a^{\alpha \beta} + \mu \left( A^{\alpha \beta} - a^{\alpha \beta} \right).
    \end{split}
\end{equation}
The compliance matrix in the current configuration is defined as
\begin{equation}\label{eq:cabgd}
    c^{\alpha \beta \gamma \delta} = \frac{\partial^{2} W}{\partial a_{\alpha \beta} \partial a_{\gamma \delta}} = 2 \frac{\partial \tau^{\alpha \beta}}{\partial a_{\gamma \delta}}
\end{equation}

Equation~\eqref{eq:kirch_NH_2} can be rewritten as

\begin{equation}\label{eq:kirch_NH_3}
    \tau^{\alpha \beta} = \frac{\Lambda}{2} \left( J^{2} - 1 \right) + \mu A^{\alpha \beta} - \mu a^{\alpha \beta} = \left( \frac{\Lambda}{2} \left( J^{2} - 1 \right) - \mu \right) a^{\alpha \beta} - \mu A^{\alpha \beta} 
\end{equation}

With~\eqref{eq:kirch_NH_3}, equation~\eqref{eq:cabgd} can be written as

\begin{equation}\label{eq:cabgd_2}
    \begin{split}
        c^{\alpha \beta \gamma \delta} &= 2 \left[ \frac{\partial }{\partial a_{\gamma \delta}} \left( \frac{\Lambda}{2} \left( J^{2} - 1 \right) - \mu \right) a^{\alpha \beta} + \left( \frac{\Lambda}{2} \left( J^{2} - 1 \right) - \mu \right) \frac{\partial a^{\alpha \beta}}{\partial a_{\gamma \delta}} \right] \\ &= 2 \left[ \frac{\Lambda}{2} \frac{\partial J^{2}}{\partial a_{\gamma \delta}} a^{\alpha \beta} + \left( \frac{\Lambda}{2} \left( J^{2} - 1 \right) - \mu \right) \frac{1}{2} m^{\alpha \beta \gamma \delta} \right] \\ &= \Lambda J^{2} a^{\alpha \beta} a^{\gamma \delta} + \left( \frac{\Lambda}{2} \left( J^{2} - 1 \right) - \mu \right) m^{\alpha \beta \gamma \delta}
    \end{split}
\end{equation}
\subsection{Discretisation of \texorpdfstring{$w_{\alpha ; \beta}$}{w-a;b}}
Using the following results,
\begin{equation}\label{eq:w_alp}
        w_{\alpha} = \boldsymbol{w} \cdot \boldsymbol{a}_{\alpha} = \left( \mathbf{N} \mathbf{w}_{e} \right)^{T} \boldsymbol{a}_{\alpha} = \mathbf{w}_{e}^{T} \mathbf{N}^{T} \boldsymbol{a}_{\alpha} = \mathbf{w}_{e}^{T} \mathbf{N}^{T} \left( \mathbf{N}_{, \alpha} \mathbf{x}_{e} \right),
    \end{equation}

    \begin{equation}\label{eq:w_n}
        w \approx \boldsymbol{w} \cdot \boldsymbol{n} = \left( \mathbf{N} \mathbf{w}_{e} \right)^{T} \boldsymbol{n} = \mathbf{w}_{e}^{T} \mathbf{N}^{T} \boldsymbol{n},
    \end{equation}
and
    \begin{equation}\label{eq:w_cov}
    w_{; \alpha} \approx \mathbf{N}_{, \alpha} \mathbf{w}_{e} + \mathbf{N} \mathbf{w}_{e , \alpha} = \mathbf{N}_{, \alpha} \mathbf{w}_{e},
    \end{equation}
the following can be calculated
    \begin{equation}\label{eq:w_alp_cov}
        \begin{split}
            w_{\alpha; \beta} &= \boldsymbol{w}_{; \beta} \boldsymbol{a}_{\alpha} + w b_{\alpha \beta} \\ &\approx \left( \mathbf{N} \mathbf{w}_{e} \right)_{; \beta} \cdot \boldsymbol{a}_{\alpha} + w \boldsymbol{n} \mathbf{N}_{, \alpha \beta} \mathbf{x}_{e} \\ &= \left( \mathbf{N}_{, \beta} \mathbf{w}_{e} \right) \cdot \boldsymbol{a}_{\alpha} + w \boldsymbol{n} \cdot \mathbf{N}_{, \alpha \beta} \mathbf{x}_{e} \\ &= \left( \mathbf{N}_{, \beta} \mathbf{w}_{e} \right) \cdot \left( \mathbf{N}_{, \alpha} \mathbf{x}_{e} \right) + w \boldsymbol{n} \mathbf{N}_{, \alpha \beta} \mathbf{x}_{e} \\ &= \mathbf{w}_{e}^{T} \mathbf{N}_{, \beta}^{T} \mathbf{N}_{, \alpha} \mathbf{x}_{e} + \left( \boldsymbol{w} \cdot \boldsymbol{n} \right) \boldsymbol{n} \cdot \mathbf{N}_{, \alpha \beta} \mathbf{x}_{e} \\ &= \mathbf{w}_{e}^{T} \mathbf{N}_{, \beta}^{T} \mathbf{N}_{, \alpha} \mathbf{x}_{e} + \mathbf{w}_{e}^{T} \mathbf{N}^{T} \boldsymbol{n} \left( \mathbf{N}_{, \alpha \beta} \cdot \boldsymbol{n} \right) \mathbf{x}_{e} \\ &= \mathbf{w}_{e}^{T} \left[ \mathbf{N}_{, \beta}^{T} \mathbf{N}_{, \alpha} + \mathbf{N}^{T} \left( \boldsymbol{n} \otimes \boldsymbol{n} \right) \mathbf{N}_{, \alpha \beta} \right] \mathbf{x}_{e}
        \end{split}
    \end{equation}

\subsection{Linearisation} \label{appendices:curvy_linearisations}
\subsubsection{Linearisation of \texorpdfstring{$\boldsymbol{a_{\alpha}}$}{a-a}}
\begin{equation}\label{eq:a_a_lin}
    \boldsymbol{a}_{\alpha} = \mathbf{N}_{,\alpha} \mathbf{x}_{e}
\end{equation}
According to the definition of the tangent vector \eqref{tangent_parametric}, the linearisation of the undiscretised
 form is
    
Linearising the parametric definition of the surface 
    \begin{equation}\label{eq:a_a_lin_3}
        \boldsymbol{x} = \boldsymbol{x} \left( \xi^{1}, \, \xi^{2} \right),
    \end{equation}
    
    \begin{equation}\label{eq:a_a_lin_4}
       \implies \Delta \boldsymbol{x} = \frac{\partial \boldsymbol{x}}{\partial \xi^{1}} \Delta \xi^{1} + \frac{\partial \boldsymbol{x}}{\partial \xi^{2}} \Delta \xi^{2},
    \end{equation}
and using the definition of the tangent vector \eqref{tangent_parametric} 
\begin{equation}\label{eq:a_a_lin_2}
    \boldsymbol{a}_{\alpha} = \boldsymbol{x}_{, \alpha} = \frac{\partial}{\partial \xi^{\alpha}} \left( \boldsymbol{x} \left( \xi^{1}, \, \xi^{2} \right) \right), 
\end{equation}
the linearisation of $\boldsymbol{a}_{\alpha}$ can be written as 
    \begin{equation}\label{eq:a_a_lin_6}
        \Delta \boldsymbol{a}_{\alpha} = \Delta \left( \boldsymbol{x}_{, \alpha}\right) = \frac{\partial}{\xi^{1}} \left[ \frac{\partial \boldsymbol{x}}{\partial \xi^{\alpha}} \right] \Delta \xi^{1} + \frac{\partial}{\partial \xi^{2}} \left[ \frac{\partial \boldsymbol{x}}{\partial \xi^{\alpha}} \right] \Delta \xi^{2}.
    \end{equation}
    
Using the property that the differential operator commutes and \eqref{eq:a_a_lin_4}     
    \begin{equation}\label{eq:a_a_lin_7}
        \Delta \boldsymbol{a}_{\alpha} = \frac{\partial}{\partial \xi^{\alpha}} \left[ \frac{\partial \boldsymbol{x}}{\partial \xi^{1}} \Delta \xi^{1} + \frac{\partial \boldsymbol{x}}{\partial \xi^{2}} \Delta \xi^{2} \right] = \frac{\partial}{\partial \xi^{\alpha}} \Delta \boldsymbol{x}.
    \end{equation}
Substituting the discretised value of $\Delta \boldsymbol{x}$ in \eqref{eq:a_a_lin_7},    
    \begin{equation}\label{eq:a_a_lin_8}
        \Delta \boldsymbol{a}_{\alpha} \approx \frac{\partial}{\partial \xi^{\alpha}} \left[ \mathbf{N} \Delta \mathbf{x}_{e} \right] = \mathbf{N}_{, \alpha} \Delta \mathbf{x}_{e}
    \end{equation}
     
\subsubsection{Linearisation of \texorpdfstring{$a_{\alpha \beta}$}{a-ab}} 
\begin{equation}\label{eq:del_aab_1}
    a_{\alpha \beta} = \boldsymbol{a}_{\alpha} \cdot \boldsymbol{a}_{\beta}
\end{equation}
Using \eqref{eq:a_a_lin_8}
\begin{equation}\label{eq:del_aab_2}
    \begin{split}
        \Delta a_{\alpha \beta} &= \Delta \left( \boldsymbol{a}_{\alpha} \cdot \boldsymbol{a}_{\beta} \right) = \Delta \boldsymbol{a}_{\alpha} \cdot \boldsymbol{a}_{\beta} + \boldsymbol{a}_{\alpha} \cdot \Delta \boldsymbol{a}_{\beta} = \boldsymbol{a}_{\beta} \cdot \Delta \boldsymbol{a}_{\alpha} + \boldsymbol{a}_{\alpha} \cdot \Delta \boldsymbol{a}_{\beta} \approx \left( \boldsymbol{a}_{\beta} \cdot \mathbf{N}_{\alpha} + \boldsymbol{a}_{\alpha} \cdot \mathbf{N}_{\beta} \right) \Delta \mathbf{x}_{e} \\ &= \left[ \left( \mathbf{N}_{, \beta} \mathbf{x}_{e} \right) \cdot \mathbf{N}_{, \alpha} \right] \Delta \mathbf{x}_{e} + \left[ \left( \mathbf{N}_{\alpha} \mathbf{x}_{e} \right) \cdot \mathbf{N}_{\beta} \right] \Delta \mathbf{x}_{e} = \left[ \mathbf{x}_{e}^{T} \mathbf{N}_{, \beta}^{T} \mathbf{N}_{\alpha} + \mathbf{x}_{e}^{T} \mathbf{N}_{, \alpha}^{T} \mathbf{N}_{, \beta} \right] \Delta \mathbf{x}_{e} \\ &= \mathbf{x}_{e}^{T} \left[ \mathbf{N}_{, \beta}^{T} \mathbf{N}_{, \alpha} + \mathbf{N}_{\alpha}^{T} \mathbf{N}_{\beta} \right] \Delta \mathbf{x}_{e}
    \end{split}
\end{equation}

\subsubsection{Linearisation of J}
Using the chain rule it can be written 
\begin{equation}\label{eq:J_lin_1}
    \Delta J = \frac{\partial J}{\partial J_{a}} \frac{\partial J_{a}}{\partial a_{\gamma \delta}} \frac{a_{\gamma \delta}}{\partial a_{\alpha}} \Delta \boldsymbol{a}_{\alpha}.
\end{equation}
Using the definitions in \eqref{area_something_man}
\begin{equation} \label{eq:J_lin_2}
    \begin{split}
        \Delta J &= \frac{1}{J_{A}} \frac{\partial \sqrt{\det a_{\gamma \delta}}}{\partial a_{\gamma \delta}} \frac{\partial a_{\gamma \delta}}{\partial \boldsymbol{a}_{\alpha}} \Delta \boldsymbol{a}_{\alpha} \\
        & = \left(\frac{1}{J_{A}}\right) \left(\frac{1}{2\,\sqrt{\det a_{\gamma \delta}}} \frac{\partial \left(\det a_{\gamma \delta}\right)}{\partial a_{\gamma \delta}}\right) \left(\frac{\partial a_{\gamma \delta}}{\partial \boldsymbol{a}_{\alpha}}\right) \Delta \boldsymbol{a}_{\alpha} \\
        & = \left(\frac{1}{J_{A}}\right) \left(\frac{1}{2 J_{a}} J_{a}^{2} a^{\gamma \delta}\right) \left(\frac{\partial a_{\gamma \delta}}{\partial \boldsymbol{a}_{\alpha}}\right) \Delta \boldsymbol{a}_{\alpha} \\
        & = \frac{J_{a}}{2 J_{A}} a^{\gamma \delta} \frac{\partial \left(\boldsymbol{a}_{\gamma} \cdot \boldsymbol{a}_{\delta}\right)}{\partial \boldsymbol{a}_{\alpha}} \Delta \boldsymbol{a}_{\alpha}\\
        & = \frac{J}{2} a^{\gamma \delta} \left[\delta^{\alpha}_{.\gamma} \boldsymbol{a}_{\delta} + \boldsymbol{a}_{\gamma} \delta_{\delta}^{.\alpha}\right] \Delta \boldsymbol{a}_{\alpha}\\
        & = \frac{J}{2} \left[a^{\alpha \delta} \boldsymbol{a}_{\delta} + a^{\alpha \gamma} \boldsymbol{a}_{\gamma}\right] \Delta \boldsymbol{a}_{\alpha}\\
        & = \frac{J}{2}\left(2 \boldsymbol{a}^{\alpha}\right) \Delta \boldsymbol{a}_{\alpha}\\
        \implies \Delta J & = J \boldsymbol{a}^{\alpha} \Delta \boldsymbol{a}_{\alpha}
    \end{split}
\end{equation}
With the discretised definitions \eqref{eq:J_lin_2} is given as
\begin{equation}\label{eq:J_lin_3}
    \Delta J \approx J \boldsymbol{a}^{\alpha} \cdot \mathbf{N}_{,\alpha} \Delta \mathbf{x}_{e}.
\end{equation}

\subsubsection{Linearisation of \texorpdfstring{$a^{\alpha \beta}$}{contravariant metric}}
Using \eqref{contravar_def}, the definition \eqref{contravar_metric_def} can be rewritten as
\begin{equation}\label{eq:del_a_ab_1}
    a^{\alpha \beta} = \frac{1}{a} e^{\alpha \gamma} a_{\gamma \delta} e^{\delta \beta}
\end{equation}
where $a$ is the determinant of the metric tensor and $e^{\alpha \gamma}$ is the unit alternator.
Thus, the linearisation of \eqref{eq:del_a_ab_1} is
\begin{equation}\label{eq:del_a_ab_2}
    \Delta a^{\alpha \beta} = \frac{\partial a^{\alpha \beta}}{\partial a_{\rho \phi}} \Delta a_{\rho \phi} = \frac{\partial \left[ a^{-1} e^{\alpha \gamma} a_{\gamma \delta} e^{\delta \beta} \right]}{\partial a_{\rho \phi}} \Delta a_{\rho \phi} = \underbrace{\frac{\partial a^{-1}}{\partial a_{\rho \phi}} e^{\alpha \gamma} a_{\gamma \delta} e^{\delta \beta} \Delta a_{\rho \phi}}_{I} + \underbrace{\frac{1}{a} e^{\alpha \gamma} \frac{\partial a_{\gamma \delta}}{\partial a_{\rho \phi}} e^{\delta \beta} \Delta a_{\rho \phi}}_{II}.
\end{equation}
Here
\begin{equation}\label{eq:del_a_ab_3}
    \begin{split}
        I &= \frac{\partial a^{-1}}{\partial a_{\rho \phi}} e^{\alpha \gamma} a_{\gamma \delta} e^{\delta \beta} \Delta a_{\rho \phi} = \left( -1 \right) a^{-2} a a^{\rho \phi} e^{\alpha \gamma} a_{\gamma \delta} e^{\delta \beta} \Delta a_{\rho \phi} \\ &= - \frac{1}{a} a^{\rho \phi} e^{\alpha \gamma} a_{\gamma \delta} e^{\delta \beta} \Delta a_{\rho \phi} = - a^{\rho \phi} \left[ \frac{1}{a} e^{\alpha \gamma} a_{\gamma \delta} e^{\delta \beta} \right] \left[ \boldsymbol{a}_{\rho} \Delta \boldsymbol{a}_{\phi} + \boldsymbol{a}_{\phi} \Delta \boldsymbol{a}_{\rho} \right] \\ &= - a^{\rho \phi} a^{\alpha \beta} \left[ \boldsymbol{a}_{\rho} \Delta \boldsymbol{a}_{\phi} + \boldsymbol{a}_{\phi} \Delta \boldsymbol{a}_{\rho} \right] = \underbrace{- a^{\rho \phi} a^{\alpha \beta} \boldsymbol{a}_{\rho} \Delta \boldsymbol{a}_{\phi}}_{\rho \rightarrow \gamma, \, \phi \rightarrow \delta} - \underbrace{a^{\rho \phi} a^{\alpha \beta} \boldsymbol{a}_{\phi} \Delta \boldsymbol{a}_{\rho}}_{\rho \rightarrow \delta, \, \phi \rightarrow \delta} \\ &= - a^{\gamma \delta} a^{\alpha \beta} \boldsymbol{a}_{\gamma} \Delta \boldsymbol{a}_{\delta} - a^{\delta \gamma} a^{\alpha \beta} \boldsymbol{a}_{\gamma} \Delta \boldsymbol{a}_{\delta} = - \left( a^{\gamma \delta} + a^{\delta \gamma} \right) a^{\alpha \beta} \boldsymbol{a}_{\gamma} \Delta \boldsymbol{a}_{\delta} \\ &= - 2 a^{\gamma \delta} a^{\alpha \beta} \boldsymbol{a}_{\gamma} \Delta \boldsymbol{a}_{\delta}.
\end{split}
\end{equation}

As the metric tensor is symmetric
\begin{equation}\label{eq:daabaab}
\frac{\partial a_{\gamma \delta}}{\partial a_{\rho \phi}} = \frac{1}{2} \frac{\partial \left[ a_{\gamma \delta}+ a_{\delta \gamma} \right]}{\partial a_{\rho \phi}},
\end{equation}

\begin{equation}\label{eq:del_a_ab_5}
    \begin{split}
        II &= \frac{1}{a} e^{\alpha \gamma} \frac{\partial a_{\gamma \delta}}{\partial a_{\rho \phi}} e^{\delta \beta} \Delta a_{\rho \phi}\\
        &= \frac{1}{a} e^{\alpha \gamma} \left[\frac{1}{2} \frac{\partial \left( a_{\gamma \delta}+ a_{\delta \gamma} \right)}{\partial a_{\rho \phi}}\right] e^{\delta \beta} \Delta a_{\rho \phi}\\    
        &= \frac{1}{a} e^{\alpha \gamma} \frac{1}{2} \frac{\partial a_{\gamma \delta}}{\partial a_{\rho \phi}} e^{\delta \beta} \Delta a_{\rho \phi} + \frac{1}{a} e^{\alpha \gamma} \frac{1}{2} \frac{\partial a_{\delta \gamma}}{\partial a_{\rho \phi}} e^{\delta \beta} \Delta a_{\rho \phi} \\
        &= \frac{1}{2a} e^{\alpha \gamma} \left[ \delta^{\rho}_{\gamma} \delta^{\phi}_{\delta} \right] e^{\delta \beta} \left[ \boldsymbol{a}_{\rho} \Delta \boldsymbol{a}_{\phi} + \boldsymbol{a}_{\phi} \Delta \boldsymbol{a}_{\rho} \right] + \frac{1}{2a} e^{\alpha \gamma} \left[ \delta^{\rho}_{\delta} \delta^{\phi}_{\gamma} \right] e^{\delta \beta} \left[ \boldsymbol{a}_{\rho} \Delta \boldsymbol{a}_{\phi} + \boldsymbol{a}_{\phi} \Delta \boldsymbol{a}_{\rho} \right] \\
        &= \frac{1}{2a} e^{\alpha \gamma} e^{\delta \beta} \left[ \boldsymbol{a}_{\rho} \delta^{\rho}_{\gamma} \Delta \boldsymbol{a}_{\phi} \delta^{\phi}_{\delta} + \boldsymbol{a}_{\phi} \delta^{\phi}_{\delta} \Delta \boldsymbol{a}_{\rho} \delta^{\rho}_{\delta} \right] + \frac{1}{2a} e^{\alpha \gamma} e^{\delta \beta} \left[ \boldsymbol{a}_{\rho} \delta^{\rho}_{\delta} \Delta \boldsymbol{a}_{\phi} \delta^{\phi}_{\gamma} + \boldsymbol{a}_{\phi} \delta^{\phi}_{\gamma} \Delta \boldsymbol{a}_{\rho} \delta^{\rho}_{\delta} \right] \\
        &= \frac{1}{2a} e^{\alpha \gamma} e^{\delta \beta} \left[ \boldsymbol{a}_{\gamma} \Delta \boldsymbol{a}_{\delta} + \boldsymbol{a}_{\delta} \Delta \boldsymbol{a}_{\gamma} \right] + \frac{1}{2a} e^{\alpha \gamma} e^{\delta \beta} \left[ \boldsymbol{a}_{\delta} \Delta \boldsymbol{a}_{\gamma} + \boldsymbol{a}_{\gamma} \Delta \boldsymbol{a}_{\delta} \right] \\
        &= \frac{1}{a} e^{\alpha \gamma} e^{\delta \beta} \left[ \boldsymbol{a}_{\gamma} \Delta \boldsymbol{a}_{\delta} + \boldsymbol{a}_{\delta} \Delta \boldsymbol{a}_{\gamma} \right] = \frac{1}{a} e^{\alpha \gamma} e^{\delta \beta} \left[ \boldsymbol{a}_{\gamma} \Delta \boldsymbol{a}_{\delta} + \boldsymbol{a}_{\delta} \Delta \boldsymbol{a}_{\gamma} \right] \\
        &= \frac{1}{a} e^{\alpha \gamma} e^{\delta \beta} \boldsymbol{a}_{\gamma} \Delta \boldsymbol{a}_{\delta} + \underbrace{\frac{1}{a} e^{\alpha \gamma} e^{\delta \beta} \boldsymbol{a}_{\delta} \Delta \boldsymbol{a}_{\gamma}}_{\gamma \rightarrow \delta, \, \delta \rightarrow \gamma} = \frac{1}{a} e^{\alpha \gamma} e^{\delta \beta} \boldsymbol{a}_{\gamma} \Delta \boldsymbol{a}_{\delta} + \frac{1}{a} e^{\alpha \delta} e^{\gamma \beta} \boldsymbol{a}_{\gamma} \Delta \boldsymbol{a}_{\delta} \\
        &= \frac{1}{a} \left( e^{\alpha \gamma} e^{\delta \beta} + e^{\alpha \delta} e^{\gamma \beta} \right) \boldsymbol{a}_{\gamma} \Delta \boldsymbol{a}_{\delta}
    \end{split}
\end{equation}
Using \eqref{eq:del_a_ab_3}, \eqref{eq:del_a_ab_5}, and the appropriate discretisations, equation \eqref{eq:del_a_ab_2} can be written as
\begin{equation}\label{eq:del_a_ab_6}
\Delta a^{\alpha \beta} = -2 a^{\alpha \beta} a^{\gamma \delta} \boldsymbol{a}_{\gamma} \Delta \boldsymbol{a}_{\delta} + \frac{1}{a} \left( a^{\alpha \gamma} e^{\delta \beta} + e^{\alpha \delta} e^{\gamma \beta} \right) \boldsymbol{a}_{\gamma} \Delta \boldsymbol{a}_{\delta} = m^{\alpha \beta \gamma \delta} \boldsymbol{a}_{\gamma} \Delta \boldsymbol{a}_{\delta} \approx m^{\alpha \beta \gamma \delta} \boldsymbol{a}_{\gamma} \mathbf{N}_{,\delta} \Delta \mathbf{x}_{e}
\end{equation}

where

\begin{equation}\label{eq:mabgd}
m^{\alpha \beta \gamma \delta} = \frac{1}{a} \left( e^{\alpha \gamma} e^{\delta \beta} + e^{\alpha \delta} e^{\gamma \beta} \right) - 2 a^{\alpha \beta} a^{\gamma \delta}
\end{equation}

\subsubsection{Linearisation of \texorpdfstring{$\tau^{\alpha \beta}$}{Kircchhof stress tensor}}
Given
\begin{equation}\label{eq:sig_ab}
    \sigma^{\alpha \beta} = \frac{\mu}{J} \left( A^{\alpha \beta} - \frac{a^{\alpha \beta}}{J^{2}} \right) 
    \end{equation}
and the definition of the Kircchhof stress tensor,     
    \begin{equation}\label{eq:kirch}
    \boldsymbol{\tau} = J \boldsymbol{\sigma}
    \end{equation}
the components of the tensor can be written as
    \begin{equation}\label{eq:kirch_2}
    \tau^{\alpha \beta} = \mu \left( A^{\alpha \beta} - a^{\alpha \beta} J^{-2} \right).
    \end{equation}
Thus the linearisation of \eqref{eq:kirch_2} is
    \begin{equation}\label{eq:del_kirch}
        \Delta \tau^{\alpha \beta} = \mu \Delta \left( A^{\alpha \beta} \right) - \mu \Delta \left( a^{\alpha \beta} J^{-2} \right) = \underbrace{\mu \Delta \left( A^{\alpha \beta} \right)}_{\substack{= 0 \\ \text{as reference metric}}} - \frac{\mu}{J^{2}} \Delta a^{\alpha \beta} - \mu a^{\alpha \beta} \Delta \left( J^{-2} \right)
    \end{equation}
    
    \begin{equation}\label{eq:del_kirch_2}
    \implies \Delta \tau^{\alpha \beta} = - \frac{\mu}{J^{2}} \Delta a^{\alpha \beta} - \mu a^{\alpha \beta} - \frac{2}{J^{3}} \Delta J.
    \end{equation}
Using~\eqref{eq:J_lin_3} and~\eqref{eq:del_a_ab_6}, it is possible to rewrite~\eqref{eq:del_kirch_2} as
\begin{equation}\label{eq:del_kirch_3}
    \begin{split}
        \Delta \tau^{\alpha \beta} &\approx - \frac{\mu}{J^{2}} m^{\alpha \beta \gamma \delta} \boldsymbol{a}_{\gamma} \mathbf{N}_{, \delta} \Delta \mathbf{x}_{e} + \frac{2}{J^{3}} \mu a^{\alpha \beta} \left( J \boldsymbol{a}^{\delta} \mathbf{N}_{\delta} \Delta \mathbf{x}_{e} \right) \\ &= - \frac{\mu}{J^{2}} \left[ 2 a^{\alpha \beta} a^{\delta \gamma} \boldsymbol{a}_{\gamma} \mathbf{N}_{, \delta} \Delta \mathbf{x}_{e} - m^{\alpha \beta \gamma \delta} \boldsymbol{a}_{\gamma} \mathbf{N}_{, \delta} \Delta \mathbf{x}_{e}\right] \\ &= \frac{\mu}{J^{2}} \left[ 2 a^{\alpha \beta} a^{\delta \gamma} - m^{\alpha \beta \gamma \delta} \right] \boldsymbol{a}_{\gamma} \mathbf{N}_{, \delta} \Delta \mathbf{x}_{e} \\ &= \frac{\mu}{J^{2}} \left[ 4 a^{\alpha \beta} a^{\delta \gamma} - \frac{1}{a} \left( a^{\alpha \gamma} e^{\delta \beta} + e^{\alpha \delta} e^{\gamma \beta}\right)  \right] \boldsymbol{a}_{\gamma} \cdot \mathbf{N}_{, \delta} \Delta \mathbf{x}_{e}
    \end{split}
\end{equation}


\end{appendices}
\end{document}